\section{Paths through a sparse set of obstacles}\label{lat:section}

The comparison of particle laws will move labels along paths of usable
cells.  Two such paths may each be long, but their mutual encounters must
have a bounded exponential moment. We combine a purely geometric skeleton
law from the companion article with the joint bad-cell bound. The detours
and their probability estimates are proved here; the bad cells need not be
independent or translation invariant.

Write $\mathcal L_K=(\Z/K\Z)^3$, with periodic maximum distance $d_\infty$.
Nearest-neighbor edges join vertices differing by one in one coordinate;
$*$-edges join distinct vertices at maximum distance one.  A random datum
$\mathcal D$ determines a bad set $\mathcal B\subset\mathcal L_K$.
All other vertices are called good.  Our assumption is
\begin{equation}\label{lat:bad-law}
 \PP(A\subset\mathcal B)\le p^{|A|}
 \qquad(A\subset\mathcal L_K).
\end{equation}
An ordered pair $(v,w)$ is admissible if the forward displacement of its
first coordinate has an integer representative $n$ with
$\lceil K/10\rceil\le n\le\lfloor K/8\rfloor$, and its two transverse
displacements have representatives $d_2,d_3$ satisfying $|d_i|\le n$.
These representatives determine the lift displacement $(n,d_2,d_3)$ uniquely.
Set $\upsilon_0=1/2000$.
For two simple paths $\pi,\pi'$ define
\begin{equation}\label{lat:contact-definition}
 J_R(\pi,\pi')=
 \sum_{z\in\pi}\sum_{z'\in\pi'}
              \ind_{\{d_\infty(z,z')\le R\}}.
\end{equation}
Here a simple path visits each vertex at most once.

\begin{proposition}[Good paths with bounded encounters]\label{lat:paths}
Fix an integer $R\ge1$.  There are constants
$p_0,\kappa>0$, $C_{\rm path}<\infty$, and an integer $K_0$, depending
only on $R$, with the following property.  Suppose $K\ge K_0$ and
\eqref{lat:bad-law} holds with $p\le p_0$.  For every admissible ordered
pair $(v,w)$ there is a $\mathcal D$-measurable event $\mathcal E_{v,w}$
of probability at least $1/2$.  On this event there is a measurable
probability kernel on simple good nearest-neighbor paths from $v$ to $w$.
Two independent draws from this kernel, conditional on $\mathcal D$,
satisfy
\begin{equation}\label{lat:encounter}
 \E\left[
   \ind_{\mathcal E_{v,w}}
   \E\left[e^{\kappa J_R(\pi,\pi')}\mid\mathcal D\right]
       \right]\le C_{\rm path}.
\end{equation}
The constants are uniform in the datum, its law, and the endpoint pair.
For every sufficiently large $K$, at least $\upsilon_0K^6$ ordered pairs
are admissible.
\end{proposition}

The event is determined before drawing the paths.  The expectation in
\eqref{lat:encounter} averages over one common environment and two fresh
path samples; it is not a bound for every fixed environment.  We will
use the proposition with $R=30$.

\subsection{The geometric skeleton input}\label{lat:skeleton-law}

The geometric input is the explicit law in
\cite[Section~8.1, pp.~22--24; Lemma~8.1 and
Equations~(24)--(25), p.~23]{SoftCompanion}.
Its use of unpredictability and intersection moments follows Benjamini,
Pemantle and Peres~\cite[Theorem~1.3 and Lemma~3.1]{BPP1998}; finite-endpoint
synchronization averages were developed by
Abbe, Massouli\'e, Montanari, Sly and Srivastava
\cite[Section~6, Lemma~6.1]{AbbeMassoulieMontanariSlySrivastava2018}.
Garban and Spencer~\cite[Theorem~2.4 and Lemma~2.5]{GarbanSpencer2023}
adapt this path-averaging method to symmetry breaking in classical
disordered spin systems.
We restate the law and both estimates independently of any gas parameters.

Fix an admissible pair, choose a lift of $v$, and give $w$ the lift
$v+(n,d_2,d_3)$. For each $j\in\{2,3\}$ use independent random variables
$U^{(j)}_{\ell,k}$, uniform on $[-1,1]$, for integers $\ell,k\ge0$. Put
\begin{equation}\label{lat:velocity-law}
 v_i^{(j)}=c_{\rm vel}\sum_{\ell\ge0}2^{-\ell/4}
             U^{(j)}_{\ell,\lfloor(i-1)/2^\ell\rfloor},
 \qquad Z_t^{(j)}=\sum_{i=1}^t v_i^{(j)},\qquad Z_0^{(j)}=0.
\end{equation}
As in the companion, fix
$c_{\rm vel}=(1-2^{-1/4})/2$, so that
$c_{\rm vel}\sum_{\ell\ge0}2^{-\ell/4}=1/2$. The series converges
absolutely for every sample. The base vertices are
\begin{equation}\label{lat:nominal}
 u_t=v+\left(t,
   \left\lfloor td_2/n+Z_{\min(t,n-t)}^{(2)}\right\rfloor,
   \left\lfloor td_3/n+Z_{\min(t,n-t)}^{(3)}\right\rfloor\right),
 \qquad 0\le t\le n.
\end{equation}
The endpoints are $v,w$. Each base increment advances the first coordinate
by one and changes each transverse coordinate by at most two. Refine it
by first making the first-coordinate step, then all required steps in
coordinate two, then those in coordinate three. This takes at most five
nearest-neighbor steps. The occurrence list consists of the initial vertex
and the vertex reached after each step, without adding an extra copy at
the join between consecutive refinements. Denote the projected walk by
$\Gamma$ and retain this lift for the detour construction.

For two independent copies $\Gamma,\Gamma'$ of this law and a fixed integer
$R_s\ge1$, define the occurrence count
\begin{equation}\label{lat:occurrence-count}
 J_{\rm occ}^{(R_s)}=
 \sum_{z\in\Gamma}\ind_{\{d_\infty(z,\Gamma')\le R_s\}}
 +\sum_{z'\in\Gamma'}\ind_{\{d_\infty(z',\Gamma)\le R_s\}}.
\end{equation}
Both sums count occurrences, including any repetitions. For a fixed $B>0$,
put
\[
 S_B(\Gamma,\Gamma')=
       \sum_{z\in\Gamma}e^{-d_\infty(z,\Gamma')/B},
\]
again counting occurrences. The companion estimates are
\begin{align}
 \E e^{\alpha_s J_{\rm occ}^{(R_s)}}&\le C_{\rm hit},
 &\alpha_s&=\frac{c_{\rm hit}}{(1+R_s)^2},
       \label{lat:imported-hit}\\
 \E e^{\xi(B)S_B(\Gamma,\Gamma')}&\le C(B),
 &\xi(B)&>0.\label{lat:imported-soft}
\end{align}
Here $c_{\rm hit}>0$ and $C_{\rm hit}<\infty$ are absolute, and the second
line's constants depend only on $B$. Both bounds are uniform in $K$ and
the admissible endpoint pair. The law and these two bounds are the sole
companion input; no interaction, density, temperature, or many-body measure
occurs in their hypotheses. We use independent skeleton seeds, also
independent of the exposed datum $\mathcal D$.

For later use we convert the occurrence count to our pair count. Every
first-coordinate level of a lifted skeleton contains at most five
occurrences: level zero has only the initial vertex, and all vertices
created in the $t$th refinement have first coordinate $v_1+t$.
The first-coordinate span is $n\le K/8$, so the periodic distance between
two such levels is their ordinary distance. A fixed occurrence of $\Gamma$
can therefore have at most $5(2R+1)$ opposing occurrences within distance
$R$. With $R_s=R$, this proves
\begin{equation}\label{lat:pair-occurrence}
 J_R(\Gamma,\Gamma')\le 5(2R+1)J_{\rm occ}^{(R)}.
\end{equation}
The left side uses the double occurrence sum from
Equation~\eqref{lat:contact-definition}. This conversion uses the actual
refinement law, as well as its moment bounds.

\subsection{Finite components and their exterior boundaries}

We record the deterministic facts used to move a lattice path around
bad vertices.  For a finite $*$-connected nonempty set $C\subset\Z^3$,
let $\operatorname{fill}(C)$ consist of $C$ and all finite
nearest-neighbor components of $\Z^3\setminus C$.  The remaining
component is the exterior.  There is only one infinite component:
outside any box containing $C$, nearest-neighbor vertices communicate,
and every infinite component meets that region.

\begin{lemma}[Exterior shell]\label{lat:boundary}
The set $\operatorname{fill}(C)$ lies in the coordinate bounding box of
$C$.  The exterior vertices at maximum distance one from $C$ form a
nearest-neighbor connected set, denoted $\partial_*^{\rm ext}C$,
with at most $26|C|$ vertices.  If $C$ is a $*$-component of a larger
bad set, every vertex of this shell is good.
\end{lemma}
\begin{proof}
A vertex outside the coordinate box has a coordinate ray to infinity
that misses $C$, so it is not in a finite complementary component.
For the shell's connectivity we apply Kesten's boundary-connectivity
theorem in the formulation of Tim\'ar~\cite[arXiv version~2, Theorem~4]{Timar2013}:
for a finite $*$-connected subset of $\Z^3$, its exterior
$*$-neighbors, accessible from infinity through its nearest-neighbor
complement, are nearest-neighbor connected.  These are exactly the
vertices in the shell just defined.  There are at most $26$ possible
$*$-neighbors per vertex of $C$.  A bad shell vertex would be
$*$-adjacent to $C$ and hence belong to the same bad component,
which is impossible for a vertex outside $C$.
\end{proof}

For an integer $a\ge1$, range-$a$ adjacency joins distinct torus
vertices whose maximum distance is at most $a$.

\begin{lemma}[Small periodic components]\label{lat:small-lifts}
Suppose $K>2a$, and a range-$a$ component on $\mathcal L_K$ has
$m$ vertices with $am<K$.  Each of its lifts to $\Z^3$ is a finite
component projecting bijectively onto it.  The coordinate span of a lift
is at most $a(m-1)$.
\end{lemma}
\begin{proof}
Choose a spanning tree, lift its root, and lift each tree edge by its
unique displacement in $[-a,a]^3$.  Distinct torus vertices have distinct
lifts.  The difference between any two lifted vertices has maximum norm
at most $a(m-1)$, by their tree path.  For a remaining edge, its assigned
endpoint difference minus its short displacement belongs to $K\Z^3$
and has maximum norm at most $am<K$.  It must therefore vanish.
All edges are consistent with these lifts.  The full inverse image is
their disjoint translates by $K\Z^3$; no edge connects two translates.
\end{proof}

We will also use a basic counting bound.  For each fixed range $a$, there
is a constant $A_a$ such that the number of connected sets of $m$ vertices
containing a specified root is at most $A_a^m$, on either $\Z^3$ or the
torus.  To see this, choose deterministically a rooted spanning tree for
each set and traverse it depth first.  The traversal has $2(m-1)$ steps,
each chosen from a fixed finite set of possible displacements, and its
visited vertices recover the set.  Increasing $A_a$ covers $m=1$.
Together with \eqref{lat:bad-law}, this gives
\begin{equation}\label{lat:large-animal}
 \PP(\text{a range-$a$ bad component has at least }m\text{ vertices})
       \le K^3(A_ap)^m.
\end{equation}
A component of size at least $m$ contains a connected subset of exactly
$m$ vertices, by growing a spanning tree one vertex at a time.

Fix $a_1=R+2$.  Reject the environment if it contains a range-$a_1$
bad component of size at least
\[
 m_K=\left\lfloor\frac{K}{10a_1}\right\rfloor.
\]
For large $K$, every remaining bad $*$-component has finite periodic
lifts, by Lemma~\ref{lat:small-lifts}.  Reject also if either endpoint
belongs to the projection of the fill of one of these components.
The event on which neither rejection occurs is $\mathcal E_{v,w}$.

If the fill of a lifted $*$-component of size $m$ contains a lift of
an endpoint, some root of the component lies within maximum distance
$m$ of that endpoint.  This follows from the bounding-box assertion
and the coordinate span in Lemma~\ref{lat:small-lifts}.  There are at
most $C m^3$ possible roots modulo $K$.  Counting connected sets and
using \eqref{lat:bad-law} thus yields
\begin{equation}\label{lat:rejection}
 \PP(\mathcal E_{v,w}^{c})
 \le K^3(A_{a_1}p)^{m_K}
       +C\sum_{m\ge1}m^3(A_1p)^m.
\end{equation}
The constant includes both endpoints.  First make $p_0$ small enough
that the series is at most $1/4$, then make $K_0$ large enough that
the first term is at most $1/4$ for $p\le p_0$ and $K\ge K_0$.
This proves $\PP(\mathcal E_{v,w})\ge1/2$.  We may decrease $p_0$
further below.

\subsection{Nearest-neighbor detours}

Draw the skeleton of Section~\ref{lat:skeleton-law} independently of the
datum. Its folding in Equation~\eqref{lat:nominal} fixes the final endpoint
without conditioning the random process on a future value. We now adapt
this nearest-neighbor walk to the bad set while keeping the endpoint
rejection event fixed.

On $\mathcal E_{v,w}$, replace the visits of this lifted walk to bad
components as follows.  If it meets a lifted bad $*$-component $C$,
take its first and last visits to $C$.  The vertices immediately before
and after these visits belong to $\partial_*^{\rm ext}C$: the path
from the first endpoint to the first such vertex misses $C$, and the
path from the last such vertex to the other endpoint misses $C$;
both endpoints are outside $\operatorname{fill}(C)$.  By
Lemma~\ref{lat:boundary}, join these two shell vertices by a simple
nearest-neighbor path in the shell, replacing the intervening part of
the walk.  The added vertices are all good.  Each replacement removes
at least one bad visit, and no replacement introduces a bad visit.
The process therefore terminates after finitely many steps.

Every component treated in this procedure was met by the original
nominal refinement: any remaining bad vertex still belongs to that
original walk.  Project the final walk to the torus and erase loops
in chronological order.  The result is a simple good nearest-neighbor
path from $v$ to $w$.  Every nonnominal vertex belongs to the shell of
a bad $*$-component met by $\Gamma$.  All choices in the refinement,
shell paths, and erasure can be made using fixed orderings.  Thus the
construction gives the measurable kernel in Proposition~\ref{lat:paths}.

We now bound the encounters introduced by this construction.  Use two
independent nominal samples $\Gamma,\Gamma'$ in the same environment,
and let $\pi,\pi'$ be their final simple paths.  Let $H$ range over
the range-$a_1$ bad components on the torus. Fix
$C_R=52(2R+1)^3$, which also exceeds $5(2R+1)$ in
Equation~\eqref{lat:pair-occurrence}. On $\mathcal E_{v,w}$,
\begin{align}\label{lat:detour-charge}
 J_R(\pi,\pi')
 &\le J_R(\Gamma,\Gamma')\notag\\
 &\quad+C_R\sum_H |H|
       \ind_{\{d_\infty(H,\Gamma)\le a_1,\,
                  d_\infty(H,\Gamma')\le a_1\}}.
\end{align}
For nominal walks, $J_R$ here counts vertex occurrences, which only
increases the bound.

To prove \eqref{lat:detour-charge}, first count encounters where both
vertices are nominal.  These are bounded by the first term.  If a
nonnominal vertex belongs to the shell of $C$ and encounters a nominal
vertex of the other walk, then $C$ is within distance $R+1\le a_1$
of that other walk, and $C$ meets its own nominal walk.  If two
nonnominal vertices belong to shells of $C,C'$, then
$d_\infty(C,C')\le R+2\le a_1$, so both components belong to one
range-$a_1$ component $H$ meeting both nominal walks.  Charge each
such encounter to one of its nonnominal vertices and to this $H$.
There are at most $26|H|$ possible added vertices associated to the
$*$-components inside $H$, even if several lifted copies were used:
projection can only decrease this number.  Each such vertex encounters
at most $(2R+1)^3$ vertices of the other simple path.  Charging from
either path costs at most an additional factor two.  This proves the
claimed bound.

\subsection{Averaging the environment and applying both moments}

Fix both complete nominal paths.  Their independence from $\mathcal D$
leaves \eqref{lat:bad-law} available for any prescribed collection of
bad vertices.  Expand the exponential of the second term in
\eqref{lat:detour-charge} as a product over the actual range-$a_1$
components.  A term corresponding to distinct components has disjoint
vertex sets.  Drop maximality, sum over all disjoint connected candidate
sets, and apply \eqref{lat:bad-law}.  Finally drop the disjointness
restriction and use $1+x\le e^x$.  The resulting bound is
\begin{align}\label{lat:environment-expansion}
 &\E_{\mathcal D}\left[
 \ind_{\mathcal E_{v,w}}
 \exp\left(\kappa C_R\sum_H|H|
   \ind_{\{d_\infty(H,\Gamma),d_\infty(H,\Gamma')\le a_1\}}\right)
                       \right]\notag\\
 &\qquad\le
 \exp\left(\sum_F
       (p e^{\kappa C_R})^{|F|}
       \ind_{\{d_\infty(F,\Gamma),d_\infty(F,\Gamma')\le a_1\}}
                         \right),
\end{align}
where $F$ ranges over all nonempty range-$a_1$ connected torus sets.
The factor $e^{\kappa C_R|F|}-1$ from the expansion was bounded above
by $e^{\kappa C_R|F|}$.  Thus only joint bad-set probabilities were
used; there is no independence assumption on the environment.

We bound the exponent by the soft-distance statistic in
Equation~\eqref{lat:imported-soft}. If a contributing set $F$ has size $m$,
choose an occurrence $z\in\Gamma$ within distance $a_1$ of a root of $F$.
There are at most $(2a_1+1)^3$ root choices for each $z$, and at most
$A_{a_1}^m$ connected sets at each root. A spanning-tree path in $F$ has
at most $m-1$ edges, each of length at most $a_1$ in the periodic maximum
metric. Since $F$ is also within $a_1$ of $\Gamma'$, the triangle inequality
gives
\[
 d_\infty(z,\Gamma')\le 2a_1+a_1(m-1)=a_1(m+1).
\]
This argument applies even to wrapping candidate sets; only the actual
components used for detours required the earlier lift and rejection test.
Absorb the root choices into a constant $A_R\ge1$. The exponent in
Equation~\eqref{lat:environment-expansion} is at most
\begin{equation}\label{lat:animal-series}
 \sum_{z\in\Gamma}\sum_{m\ge1}q^m
       \ind_{\{d_\infty(z,\Gamma')\le a_1(m+1)\}},
 \qquad q=A_Rp e^{\kappa C_R}.
\end{equation}
All constants here depend only on the fixed $R$. Counting the same set
more than once only increases this nonnegative upper bound.

Require $q\le e^{-2}$. If $d\le a_1(m+1)$, split $q^m$ into two equal
powers and use
\[
 q^m=q^{m/2}q^{m/2}\le q^{m/2}e^{-m}
                       \le e\,q^{m/2}e^{-d/a_1}.
\]
Summing the first factor over $m\ge1$ proves
\begin{equation}\label{lat:soft-animal-sum}
 \sum_{m\ge1}q^m\ind_{\{d\le a_1(m+1)\}}
 \le\eta(p,\kappa)e^{-d/a_1},
 \qquad \eta(p,\kappa)=\frac{e\sqrt q}{1-\sqrt q}.
\end{equation}
In particular $\eta(p,\kappa)\to0$ as $p\to0$ at fixed $\kappa$.
Set $B=a_1$. Equations~\eqref{lat:detour-charge},
\eqref{lat:pair-occurrence}, and
\eqref{lat:environment-expansion}--\eqref{lat:soft-animal-sum} give
\begin{equation}\label{lat:moment-reduction}
 \E_{\mathcal D}\left[
  \ind_{\mathcal E_{v,w}}e^{\kappa J_R(\pi,\pi')}
       \mid\Gamma,\Gamma'\right]
 \le\exp\left(\kappa C_RJ_{\rm occ}^{(R)}+
                      \eta(p,\kappa)S_B(\Gamma,\Gamma')\right).
\end{equation}
On the left, the adapted paths are used only on $\mathcal E_{v,w}$.
The environment was averaged once, with both entire skeletons fixed.

First choose $\kappa>0$ with $2\kappa C_R\le\alpha_R$, where
$\alpha_R=c_{\rm hit}/(1+R)^2$. Next choose $p_0>0$ so small that
$q\le e^{-2}$ and $2\eta(p,\kappa)\le\xi(B)$ for $p\le p_0$, as well
as the earlier endpoint rejection requirements. The constants $K_0$ can
then be enlarged to cover those rejection estimates. Cauchy--Schwarz and
both imported moments imply
\begin{align*}
 \E\exp\left(\kappa C_RJ_{\rm occ}^{(R)}+\eta S_B\right)
 &\le
 \left(\E e^{2\kappa C_RJ_{\rm occ}^{(R)}}\right)^{1/2}
 \left(\E e^{2\eta S_B}\right)^{1/2}\\
 &\le\sqrt{C_{\rm hit}C(B)}=:C_{\rm path}<\infty.
\end{align*}
Together with Equation~\eqref{lat:moment-reduction}, this proves the
annealed bound~\eqref{lat:encounter}. Both path draws use one common
datum and independent skeleton seeds. The event retains its original
mass; there is no renormalization by conditioning on acceptance and no
assertion of a bound for each fixed environment.

Finally, for $K\ge80$ the interval
$[\lceil K/10\rceil,\lfloor K/8\rfloor]$ contains at least
$K/40-1\ge K/80$ integers. For each such $n$, there are
$(2n+1)^2\ge K^2/25$ transverse choices. They are distinct modulo $K$,
since $2n<K$. Multiplying by the $K^3$ choices of starting vertex gives
at least $K^6/2000=\upsilon_0K^6$ admissible ordered pairs. This completes
Proposition~\ref{lat:paths}.
